The woman in today's Gospel puts leaven into a large amount of flour. Jesus says that she hides it there. Once mixed in, that small amount is surrounded by what it will change. Its work takes place within the dough, until the whole has been leavened. Beside this household scene, Jesus places a seed whose small beginning gives little indication of the branches in which birds will eventually find a home.

These two parables come after others that have described serious obstacles. Some of the sower's seed fails to grow. An enemy scatters weeds among the wheat. Saint John Chrysostom explains that the disciples could become discouraged by what they have heard. They are few, and Jesus has told them about resistance to the word they will preach. The mustard seed answers their fear about what such a beginning can produce. Christ knows their weakness when he entrusts his teaching to them.

Chrysostom then follows the woman's action closely. Leaven works through contact with the flour. That is why being surrounded by people who have not received the Gospel need not make the disciples' work hopeless. Christ can make his word fruitful through their presence among those people. The disciples will encounter opposition there. They will also encounter the people to whom they have been sent. Their task requires a faith they can live among others, including those who find it difficult to receive.

Paul describes the beginning of just such a life in today's Epistle. The Thessalonians received the word amid much affliction. Their conversion did not wait for a peaceful interval in which everybody welcomed their faith. They turned from idols to serve the living and true God while that decision was costly. Then something happened to their relation with other believers. People who had first followed Paul and the Lord became an example themselves. Their faith was spoken of beyond their own city, in Macedonia and Achaia.

That change from following an example to becoming one helps us understand the leaven. The people receiving the word begin to affect other people through the life it produces in them. Paul names work, charity and endurance. His thanksgiving concerns acts that other people can recognize. Saint Thomas Aquinas explains that faith must become action and continue when action becomes difficult. He sees hope sustaining patience through adversity. The Thessalonians have something to await from Christ, and that expectation gives them a reason to persevere in what they have begun.

Paul also says that they received the word with joy from the Holy Spirit. He places that joy within their affliction. Chrysostom explains the difference between the injury they suffer and the help the Spirit gives them in belonging to Christ. The hurt remains real. A person can grieve over what has happened and still be grateful to the Lord who sustains him. You do not have to appear untroubled before your faith can help someone else. The Thessalonians encouraged other believers while they themselves still needed endurance.

This makes their example possible to understand from within our own weakness. We may think that we could serve God more faithfully if the resistance around us first disappeared. We can also postpone a duty because we cannot see what good it will do. Yet Paul gives thanks for a community whose love has become active under pressure. Their neighbors met believers whose allegiance to God had changed how they acted.

Chrysostom knew that his own congregation could admire the apostles and then excuse themselves from following them. They could point to the apostles' miracles as a reason for doing nothing themselves. He therefore pointed to the apostles' freedom from greed and their refusal to live for praise. These virtues were within his listeners' reach. A striking gift could coexist with wrongdoing, whereas a life directed by charity gave people reason to attend to the teaching.

The Collect asks God for that connection between what we understand and what we do. It includes our words as well as our deeds. What we have considered in prayer is to affect what another person hears from us. If we understand Christ's mercy and then speak with contempt, our speech contradicts the good we have recognized. The person hearing us receives the contempt. He cannot receive the charitable intention we never put into words.

The prayer at the offering asks for more help still. The Secret asks that the sacrifice cleanse and renew us, and that its effects include direction and protection. Cleansing concerns the wrong that has entered our conduct. Renewal concerns the life God restores in us. We also need his direction for the choices that remain ahead, and his protection if the good already received is to endure. These petitions concern the same people whose lives can become an example.

There is room in this Mass for an urgent admission of need. Before the Gospel, and again at the Offertory, the psalmist cries to God from the depths. Having heard the promise of growth, we still ask to be heard. We can acknowledge a failure without deciding that Christ can do nothing further through us. The Thessalonians' example gives us reason to ask for the grace by which our own conduct can be made faithful. Paul begins by thanking God for them because their conversion and perseverance are gifts for which God deserves thanksgiving.

After Communion, the Church asks that those who have received heavenly food continue to desire it. Being nourished leads to a request for continuing hunger. We need the life Christ gives as we return to the words and deeds for which we are responsible. The desire to receive his gifts again belongs to perseverance. We have more to learn, more need of cleansing, and further occasions to practice charity.

If someone's hostility has made you want to answer with contempt, speak to that person patiently the next time you meet. You can tell the truth without seeking to wound. That response allows the faith you profess to affect the encounter itself. You cannot decide how the other person will receive it. You can seek God's help to make your own words consistent with the charity he asks of you.

Paul ends his thanksgiving by describing the Son whom the Thessalonians await from heaven. God has raised Jesus from the dead. Saint Thomas explains their expectation in terms of sharing his resurrection and receiving deliverance from sin and its punishment. Their example has already strengthened other believers, and they still await resurrection and deliverance with the risen Lord.

You belong to a people who can give thanks for grace already received and still ask for help in what lies ahead. Christ has given you his word, and in this Mass you seek the nourishment by which to live it. He can make your faith fruitful while your need for him remains. As you await his coming, the patience and charity he forms in you can strengthen people who are also learning to trust him. The risen Lord whom we await will bring to fulfillment the life for which we now depend on his gifts.
